\documentclass[10pt,a4paper]{amsart}
\usepackage[margin=19mm]{geometry}
\usepackage{amsmath,amssymb,mathtools}
\usepackage[scr=rsfs,cal=euler]{mathalfa}
\usepackage{xfrac}
\usepackage[unicode,hidelinks,bookmarks=false]{hyperref}
\newcommand{\ie}{i.e.\ }
\newcommand{\cf}{cf.\ }
\newcommand{\resp}{respectively\ }
\newcommand{\Subsection}{\subsection}
\providecommand{\nobreakdash}{\nobreak\hbox{-}}
\setlength{\emergencystretch}{2em}
\multlinegap0pt
\usepackage{amssymb}
\usepackage{bm}
\usepackage{float}
\usepackage{dsfont}
\newtheorem{assumption}{Assumption}
\newtheorem{lemma}{Lemma}
\newcommand{\E}{\mathrm{E}}
\title{\bf Semiparametric Off-Policy Inference for Optimal Policy Values under Possible Non-Uniqueness}
\author{Haoyu Wei}
\date{Source: arXiv:2505.13809v5 (CC BY 4.0)}
\begin{document}
\maketitle

\noindent\emph{Source and licence.} \bf Semiparametric Off-Policy Inference for Optimal Policy Values under Possible Non-Uniqueness. Version arXiv:2505.13809v5 (CC BY 4.0), distributed by arXiv under the Creative Commons Attribution 4.0 International licence, http://creativecommons.org/licenses/by/4.0/, as recorded in the arXiv licence metadata for this exact version and shown on the abstract page https://arxiv.org/abs/2505.13809v5. Full licence notice: \texttt{licenses/arxiv-2505.13809-CC-BY-4.0.txt}.\par\smallskip

\noindent\textit{Editorial scope.} Counted unit: the article's lemma lem:finite\_mdp\_uniform\_expansion (source0.tex lines 3168--3213). The article's complete original proof of it is delivered with it (source0.tex lines 3215--3491, one \texttt{proof} environment of 277 lines). No mathematics is written, repaired or re-derived: the statement and the proof are the article's own lines, in the article's own notation, and no retained line is substituted or re-flowed.  Retained uncounted as the source-local context the proof needs: lines 581--609; lines 2743--2774.  Nothing was omitted from any retained span, so the package carries no omission record.  No retained line contains a citation, so the wrapper carries no bibliography and no bibliography entry.  No retained line contains a figure environment, an image-inclusion command or drawing code, so no image is delivered and the package declares no asset dependency.  Editorial formatting only, in addition: an AMS \texttt{amsart} preamble that restates the article's own declarations and macros used by the retained lines, namely the environments \texttt{assumption}, \texttt{lemma} on separate counters, together with the standard packages those lines need.  The retained environments print in this document as Assumption 1, Lemma 1 in order of appearance, whereas the article prints the same items with the numbering of its own full document.  The complete native source of the article as distributed in the arXiv e-print for this exact version is bundled unchanged as \texttt{sources/178401/source0.tex}.
\begin{assumption}[Finite-MDP Local Regularity in
$\mathcal M^\dagger$]\label{ass:regularity}
    The state and action spaces are finite,
    $|R|\leq\overline c_R<\infty$, $\gamma<1$, and $\Pi$ is a closed
    class of stationary
    Markov policies.  There is a constant $c>0$ such that the
    stationary behavior state law satisfies $f_b(s)\geq c$ for every
    state reachable under a policy in $\Pi$, and
    $b(a\mid s)\geq c$ whenever $\pi(a\mid s)>0$ for some
    $\pi\in\Pi$.
    For every regular parametric submodel
    $\{P_\epsilon:\epsilon\in(-\epsilon_0,\epsilon_0)\}$ through $P_0$,
    the initial law $\mu_\epsilon$, conditional reward mean
    $r_\epsilon(s,a)$, and transition kernel
    $p_\epsilon(s'\mid s,a)$ satisfy
    \[
        \begin{aligned}
        \|\mu_\epsilon-\mu_0-\epsilon\dot\mu\|_1
        =o(|\epsilon|),\qquad
        \max_{s,a}|r_\epsilon(s,a)-r_0(s,a)
        -\epsilon\dot r(s,a)|
        & =o(|\epsilon|),\\
        \text{and} \qquad \max_{s,a}\sum_{s'}
        |p_\epsilon(s'\mid s,a)-p_0(s'\mid s,a)
        -\epsilon\dot p(s'\mid s,a)|
        &=o(|\epsilon|).
        \end{aligned}
    \]
\end{assumption}

\begin{lemma}[One-Sided Occupancy Stability]
\label{lem:diff_polcies_upper}
    For two policies under the same law $P$,
    \[
        \|d_P^{\pi_2}-d_P^{\pi_1}\|_{\mathrm{TV}}
        \leq
        \frac{\gamma}{1-\gamma}
        \E_{S\sim d_P^{\pi_1}}
        \!\left[\operatorname{TV}
        \{\pi_2(\cdot\mid S),\pi_1(\cdot\mid S)\}\right].
    \]
    Equivalently, when the occupancy measures admit densities,
    \[
        \|\rho_P^{\pi_2}-\rho_P^{\pi_1}\|_{L_1(f_b)}
        \leq
        \frac{2\gamma}{1-\gamma}
        \E_{S\sim d_P^{\pi_1}}
        \!\left[\operatorname{TV}
        \{\pi_2(\cdot\mid S),\pi_1(\cdot\mid S)\}\right].
    \]
    If $\nu_P^\pi(\mathrm ds,\mathrm da)
    :=d_P^\pi(\mathrm ds)\pi(\mathrm da\mid s)$ denotes the discounted
    state--action occupancy measure, then
    \[
        \|\nu_P^{\pi_2}-\nu_P^{\pi_1}\|_{\mathrm{TV}}
        \leq
        \frac{1}{1-\gamma}
        \E_{S\sim d_P^{\pi_1}}
        \!\left[\operatorname{TV}
        \{\pi_2(\cdot\mid S),\pi_1(\cdot\mid S)\}\right].
    \]
\end{lemma}

\begin{lemma}[Uniform Fixed-Policy Expansion in the Finite MDP]
\label{lem:finite_mdp_uniform_expansion}
    Under Assumption~\ref{ass:regularity}, define
    \[
        r_{\epsilon,\pi}(s)
        :=\sum_a\pi(a\mid s)r_\epsilon(s,a),
        \qquad
        K_{\epsilon,\pi}(s,s')
        :=\sum_a\pi(a\mid s)p_\epsilon(s'\mid s,a),
    \]
    with
    $\dot r_\pi(s):=\sum_a\pi(a\mid s)\dot r(s,a)$ and
    $\dot K_\pi(s,s')
      :=\sum_a\pi(a\mid s)\dot p(s'\mid s,a)$,
    and $G_{0,\pi}:=(I-\gamma K_{0,\pi})^{-1}$.  Then
    \[
        V_{P_\epsilon}^{\pi}
        =V_{P_0}^{\pi}
        +\epsilon G_{0,\pi}
        \{\dot r_\pi+\gamma\dot K_\pi V_{P_0}^{\pi}\}
        +o(|\epsilon|)
    \]
    uniformly over $\pi\in\Pi$, where the remainder is in sup norm.
    Consequently,
    \begin{equation}\label{eq:finite_mdp_uniform_expansion}
        \sup_{\pi\in\Pi}
        \left|
        \Psi(P_\epsilon;\pi)-\Psi(P_0;\pi)
        -\epsilon\dot\Psi_{P_0}^{\pi}
        \right|
        =o(|\epsilon|),
    \end{equation}
    with
    \[
        \dot\Psi_{P_0}^{\pi}
        =
        \dot\mu^\top V_{P_0}^{\pi}
        +
        \mu_0^\top G_{0,\pi}
        \{\dot r_\pi+\gamma\dot K_\pi V_{P_0}^{\pi}\}.
    \]
    In the auxiliary experiment $\mathcal M^\dagger$,
    $\dot\Psi_{P_0}^{\pi}
      =\E_{P_0^\dagger}
      [D_{P_0}^{\dagger,\pi}\dot\ell^\dagger]$.
\end{lemma}

\begin{proof}
    Throughout this proof, the matrix norm is the maximum absolute
    row-sum norm, and the vector norm is the sup norm.  Because
    $K_{\epsilon,\pi}$ is stochastic,
    \[
        \sup_{\epsilon,\pi}
        \|(I-\gamma K_{\epsilon,\pi})^{-1}\|_\infty
        \leq
        \sum_{m=0}^{\infty}\gamma^m
        \|K_{\epsilon,\pi}^m\|_\infty
        =\frac1{1-\gamma}.
    \]
    Write
    $G_{\epsilon,\pi}:=(I-\gamma K_{\epsilon,\pi})^{-1}$,
    $\Delta r_{\epsilon,\pi}:=
      r_{\epsilon,\pi}-r_{0,\pi}$, and
    $\Delta K_{\epsilon,\pi}:=
      K_{\epsilon,\pi}-K_{0,\pi}$.
    The componentwise expansions in
    Assumption~\ref{ass:regularity} remain uniform after averaging
    with any policy, because all policy weights are nonnegative and
    sum to one.  More explicitly,
    \[
        \begin{aligned}
        \sup_{\pi}\|
        \Delta r_{\epsilon,\pi}-\epsilon\dot r_\pi
        \|_\infty
        &\leq
        \max_{s,a}|r_\epsilon(s,a)-r_0(s,a)
        -\epsilon\dot r(s,a)|
        =o(|\epsilon|),\\
        \sup_{\pi}\|
        \Delta K_{\epsilon,\pi}-\epsilon\dot K_\pi
        \|_\infty
        &\leq
        \max_{s,a}\sum_{s'}
        |p_\epsilon(s'\mid s,a)-p_0(s'\mid s,a)
        -\epsilon\dot p(s'\mid s,a)|
        =o(|\epsilon|).
        \end{aligned}
    \]
    In particular,
    $\sup_\pi\|\Delta r_{\epsilon,\pi}\|_\infty=O(|\epsilon|)$
    and
    $\sup_\pi\|\Delta K_{\epsilon,\pi}\|_\infty=O(|\epsilon|)$.

    The resolvent identity follows by multiplying out:
    \[
        \begin{aligned}
        G_{\epsilon,\pi}-G_{0,\pi}
        &=
        G_{\epsilon,\pi}
        \{(I-\gamma K_{0,\pi})
          -(I-\gamma K_{\epsilon,\pi})\}
        G_{0,\pi}\\
        &=
        \gamma G_{\epsilon,\pi}
        (K_{\epsilon,\pi}-K_{0,\pi})G_{0,\pi}
        =
        \gamma G_{\epsilon,\pi}
        \Delta K_{\epsilon,\pi}G_{0,\pi}.
        \end{aligned}
    \]
    Consequently,
    \[
        \sup_\pi\|G_{\epsilon,\pi}-G_{0,\pi}\|_\infty
        \leq
        \frac{\gamma}{(1-\gamma)^2}
        \sup_\pi\|\Delta K_{\epsilon,\pi}\|_\infty
        =O(|\epsilon|).
    \]

    We now derive the value expansion directly from the two Bellman
    equations.  Since
    $(I-\gamma K_{\epsilon,\pi})V_{P_\epsilon}^\pi
      =r_{\epsilon,\pi}$ and
    $(I-\gamma K_{0,\pi})V_{P_0}^\pi=r_{0,\pi}$,
    \[
        \begin{aligned}
        &(I-\gamma K_{\epsilon,\pi})
        (V_{P_\epsilon}^\pi-V_{P_0}^\pi)\\
        &\quad=
        r_{\epsilon,\pi}
        -(I-\gamma K_{\epsilon,\pi})V_{P_0}^\pi\\
        &\quad=
        \Delta r_{\epsilon,\pi}
        +\gamma\Delta K_{\epsilon,\pi}V_{P_0}^\pi.
        \end{aligned}
    \]
    Therefore the following identity is exact:
    \begin{equation}\label{eq:exact_value_perturbation}
        V_{P_\epsilon}^\pi-V_{P_0}^\pi
        =
        G_{\epsilon,\pi}
        \{\Delta r_{\epsilon,\pi}
        +\gamma\Delta K_{\epsilon,\pi}V_{P_0}^\pi\}.
    \end{equation}
    Bounded rewards and the resolvent bound imply
    $\sup_\pi\|V_{P_0}^\pi\|_\infty
      \leq\overline c_R/(1-\gamma)$.
    Define
    \[
        h_\pi:=\dot r_\pi+\gamma\dot K_\pi V_{P_0}^\pi.
    \]
    The primitive expansions and the uniform bound on $V_{P_0}^\pi$
    give
    \[
        \sup_\pi
        \|\Delta r_{\epsilon,\pi}
        +\gamma\Delta K_{\epsilon,\pi}V_{P_0}^\pi
        -\epsilon h_\pi\|_\infty=o(|\epsilon|),
        \qquad
        \sup_\pi\|h_\pi\|_\infty<\infty.
    \]
    Subtracting $\epsilon G_{0,\pi}h_\pi$ from
    \eqref{eq:exact_value_perturbation} yields
    \[
        \begin{aligned}
        &V_{P_\epsilon}^\pi-V_{P_0}^\pi
        -\epsilon G_{0,\pi}h_\pi\\
        &\quad=
        \epsilon(G_{\epsilon,\pi}-G_{0,\pi})h_\pi\\
        &\qquad+
        G_{\epsilon,\pi}
        \{\Delta r_{\epsilon,\pi}
        +\gamma\Delta K_{\epsilon,\pi}V_{P_0}^\pi
        -\epsilon h_\pi\}.
        \end{aligned}
    \]
    The first term is uniformly $O(\epsilon^2)$ and the second is
    uniformly $o(|\epsilon|)$.  This proves the asserted uniform
    expansion of $V^\pi$.

    Next expand the target value without suppressing the cross term:
    \[
        \begin{aligned}
        \Psi(P_\epsilon;\pi)-\Psi(P_0;\pi)
        &=(\mu_\epsilon-\mu_0)^\top V_{P_0}^\pi\\
        &\quad+\mu_0^\top
        (V_{P_\epsilon}^\pi-V_{P_0}^\pi)\\
        &\quad+(\mu_\epsilon-\mu_0)^\top
        (V_{P_\epsilon}^\pi-V_{P_0}^\pi).
        \end{aligned}
    \]
    The first two lines equal
    \[
        \epsilon\dot\mu^\top V_{P_0}^\pi
        +\epsilon\mu_0^\top G_{0,\pi}h_\pi
        +o(|\epsilon|)
    \]
    uniformly in $\pi$.  The final cross term is
    $O(\epsilon^2)$ uniformly, because both factors are
    $O(|\epsilon|)$.  This proves
    \eqref{eq:finite_mdp_uniform_expansion} and the displayed formula
    for $\dot\Psi_{P_0}^\pi$.

    It remains to identify the derivative in
    $\mathcal M^\dagger$.  Its product-factorized density can be
    written as
    \[
        p^\dagger(s_0,s,a,r,s')
        =\mu(s_0)g(s)b(a\mid s)c(r,s'\mid a,s),
    \]
    so a regular score decomposes as
    \[
        \dot\ell^\dagger
        =\dot\ell_\mu(S_0)+\dot\ell_g(S)
        +\dot\ell_b(A\mid S)
        +\dot\ell_c(R,S'\mid A,S),
    \]
    where $g$ is the transition-sampling state marginal.  The
    conditional score components satisfy
    $\E(\dot\ell_\mu)=0$,
    $\E(\dot\ell_g)=0$,
    $\E(\dot\ell_b\mid S)=0$, and
    $\E(\dot\ell_c\mid S,A)=0$.
    Differentiating the primitive laws gives
    \[
        \begin{aligned}
        \dot\mu^\top V_{P_0}^\pi
        &=\E\!\left[
        \{V_{P_0}^\pi(S_0)-\Psi(P_0;\pi)\}
        \dot\ell_\mu(S_0)\right],\\
        \dot r(s,a)
        &=\E\!\left[
        \{R-r_0(s,a)\}\dot\ell_c(R,S'\mid a,s)
        \mid S=s,A=a\right],\\
        \sum_{s'}\dot p(s'\mid s,a)V_{P_0}^\pi(s')
        &=\E\!\left[
        \{V_{P_0}^\pi(S')
        -P_0V_{P_0}^\pi(s,a)\}
        \dot\ell_c(R,S'\mid a,s)
        \mid S=s,A=a\right].
        \end{aligned}
    \]
    The centering terms in the last two lines may be omitted because
    the conditional score has mean zero.  By the Bellman equation,
    $Q_{P_0}^\pi(s,a)
      =r_0(s,a)+\gamma P_0V_{P_0}^\pi(s,a)$, and hence
    \[
        h_\pi(s)
        =
        \sum_a\pi(a\mid s)
        \E\!\left[
        \{R+\gamma V_{P_0}^\pi(S')
        -Q_{P_0}^\pi(A,S)\}\dot\ell_c
        \mid S=s,A=a\right].
    \]

    The discounted occupancy identity follows from the Neumann
    series:
    \[
        \begin{aligned}
        \mu_0^\top G_{0,\pi}h
        &=\sum_{t=0}^{\infty}\gamma^t
          \mu_0^\top K_{0,\pi}^t h\\
        &=\frac1{1-\gamma}
          \E_{S\sim d_{P_0}^\pi}\{h(S)\}.
        \end{aligned}
    \]
    Changing measure from the discounted target occupancy to the
    behavior transition marginal uses
    \[
        \omega^\pi(a,s)
        =\frac{d_{P_0}^\pi(s)\pi(a\mid s)}
        {f_b(s)b(a\mid s)}.
    \]
    Thus the second term of $\dot\Psi_{P_0}^\pi$ is
    \[
        \frac1{1-\gamma}
        \E_{P_{0,b}^{\mathrm{tr}}}\!\left[
        \omega^\pi(A,S)
        \{R+\gamma V_{P_0}^\pi(S')
        -Q_{P_0}^\pi(A,S)\}
        \dot\ell_c(R,S'\mid A,S)\right].
    \]
    Let
    $\varepsilon^\pi
      :=R+\gamma V_{P_0}^\pi(S')-Q_{P_0}^\pi(A,S)$.
    The Bellman equation gives
    $\E(\varepsilon^\pi\mid S,A)=0$.  Consequently,
    \[
        \E\{\omega^\pi\varepsilon^\pi\dot\ell_g(S)\}=0,
        \qquad
        \E\{\omega^\pi\varepsilon^\pi
        \dot\ell_b(A\mid S)\}=0.
    \]
    The initial-state component is centered and, in the product
    auxiliary experiment, is orthogonal to the transition-sampling
    score components.  Combining these orthogonality relations with
    the preceding derivative calculation gives
    \[
        \dot\Psi_{P_0}^{\pi}
        =\E_{P_0^\dagger}\!\left[
        D_{P_0}^{\dagger,\pi}(O^\dagger)
        \dot\ell^\dagger(O^\dagger)\right].
    \]
    Each summand of $D_{P_0}^{\dagger,\pi}$ is mean zero and belongs
    to the corresponding component of the full nonparametric product
    tangent space.  Therefore the displayed gradient belongs to
    $\mathcal T^\dagger$ and is the canonical gradient.

    Finally, in the finite MDP,
    $\pi\mapsto K_{0,\pi}$ and $\pi\mapsto r_{0,\pi}$ are continuous.
    The resolvent identity then makes
    $\pi\mapsto V_{P_0}^\pi$ and
    $\pi\mapsto Q_{P_0}^\pi$ continuous in sup norm.
    Lemma~\ref{lem:diff_polcies_upper}, the finite-state lower overlap
    bounds, and
    $\omega^\pi(a,s)
      =d^\pi(s)\pi(a\mid s)/\{f_b(s)b(a\mid s)\}$ make
    $\pi\mapsto\omega^\pi$ continuous in $L_1$ and hence in $L_2$
    on this finite, uniformly bounded space.  All factors in
    $D_{P_0}^{\dagger,\pi}$ are uniformly bounded, so
    $\pi\mapsto D_{P_0}^{\dagger,\pi}$ is continuous in
    $L_2(P_0^\dagger)$.
\end{proof}

\end{document}
